% META: {"id": "1NSI-WEB-CO-003", "chapitre": "1NSI-WEB-IHM", "type_objet": "corrige", "exercice_ref": "1NSI-WEB-EX-003", "capacites": ["P-IHM-03A", "P-IHM-03B"], "sources_inspiration": [], "status": "generated"}
% BEGIN-VERIFY
% cookies_client = {}
% def serveur_definit_cookie(nom, valeur):
%     cookies_client[nom] = valeur
% def client_envoie_requete_avec_cookies(url):
%     return {"url": url, "cookies": dict(cookies_client)}
% r1 = client_envoie_requete_avec_cookies("/accueil")
% assert r1 == {"url": "/accueil", "cookies": {}}
% serveur_definit_cookie("panier", "2 articles")
% r2 = client_envoie_requete_avec_cookies("/paiement")
% assert r2 == {"url": "/paiement", "cookies": {"panier": "2 articles"}}
% END-VERIFY
\begin{corrige}{1NSI-WEB-EX-003}
\textbf{1.} La requête pour \lstinline{/accueil} ne contient \textbf{aucun} cookie :
\lstinline{{"url": "/accueil", "cookies": {}}}.

\textbf{2.} La requête pour \lstinline{/paiement} contient le cookie \lstinline{panier} :
\lstinline{{"url": "/paiement", "cookies": {"panier": "2 articles"}}}.

\textbf{3.} Non : le serveur n'a pas besoin de conserver lui-même l'état du panier de
chaque visiteur. C'est le \textbf{navigateur} qui mémorise le cookie et le
\textbf{retransmet automatiquement} à chaque requête vers ce site : le serveur retrouve
ainsi l'information sans avoir eu à la stocker de son côté.
\end{corrige}
